\documentclass[10pt,a4paper]{article}
\usepackage[margin=2cm]{geometry}
\usepackage{fancyvrb,booktabs,enumitem,amsmath}
\fvset{fontsize=\scriptsize,frame=single}
\title{Lab Report -- Bayesian Networks and Autoregressive Language Models}
\author{Apurva Verma}\date{}
\begin{document}\maketitle

\section*{Part I--II: Questions 1 and 2}
\textbf{Q1.} The chain rule turns the joint probability of a whole sentence into a product of next-token conditionals, so a model only has to learn ``what comes next given what came before''. Generation then reduces to repeating one operation: sample $X_t\sim P(X_t\mid X_1,\dots,X_{t-1})$, append, repeat; the same factors also give the probability of any sentence.

\textbf{Q2.} The first-order chain $X_1\to X_2\to\dots\to X_T$ assumes the Markov property $P(X_t\mid X_1,\dots,X_{t-1})=P(X_t\mid X_{t-1})$, i.e.\ $X_t\perp (X_1,\dots,X_{t-2})\mid X_{t-1}$: given the previous word, earlier words add no information.

\section*{Part III--IV: Dataset and CPT (Question 3)}
Six lower-cased sentences, each wrapped as \texttt{<START> \dots\ <END>}; vocabulary of 10 words plus \texttt{<END>} ($|V|=11$). $P(w_j\mid w_i)=C(w_i,w_j)/\sum_k C(w_i,w_k)$.
\begin{center}\begin{tabular}{ll}\toprule Context & Non-zero probabilities\\\midrule
\texttt{<START>} & the 1.0\\
the & cat 0.250, dog 0.250, mat 0.167, rug 0.167, park 0.167\\
cat & sat 0.667, ran 0.333\\
dog & sat 0.667, ran 0.333\\
sat & on 1.0\\
ran & to 1.0\\
on, to & the 1.0\\
mat, rug, park & \texttt{<END>} 1.0\\\bottomrule\end{tabular}\end{center}
Zero-probability transitions (examples): $P(\text{cat}\mid\text{cat})=0$, $P(\text{sat}\mid\text{the})=0$, $P(\text{ran}\mid\text{sat})=0$, $P(\texttt{<END>}\mid\text{the})=0$, $P(\text{the}\mid\text{ran})=0$ (ran must be followed by to). In total the first-order model has 17 non-zero entries out of $11\times11=121$ possible transitions, so 104 transitions have probability zero.

\section*{Part V--VI: LLM prompt and inspection (Questions 4--7)}
Prompt: see appendix (the lab's suggested prompt plus the dataset and extra requirements). \texttt{ngram\_lm.py} was generated by an LLM assistant (Claude).
\textbf{Q4.} Transition counts are in \texttt{self.counts}, a \texttt{defaultdict(Counter)} keyed by the context tuple, filled in \texttt{\_\_init\_\_} while scanning \texttt{<START> \dots\ <END>}-padded sentences.
\textbf{Q5.} $P(X_t\mid X_{t-1})$ is computed in \texttt{\_\_init\_\_} in the dictionary comprehension that builds \texttt{self.probs} (count divided by the row total) and read via \texttt{dist()}.
\textbf{Q6.} \texttt{sample()} draws from the distribution with \texttt{random.choices(weights=\dots)} (option 2); \texttt{argmax()} is the greedy option 1, used only in greedy mode. Greedy always picks the single most probable word and is deterministic, sampling picks each word with its probability, so less likely words also occur.
\textbf{Q7.} For an unseen context, \texttt{dist()} returns an empty dictionary, \texttt{sample()}/\texttt{argmax()} return \texttt{None} and \texttt{generate()} stops the sentence (the \texttt{banana} test returns \texttt{None}). Without this guard a \texttt{KeyError} would be raised.

\section*{Part VII: Normalisation test (Question 8)}
For every context of both models, $\sum_v P(v\mid w)$ is $1.000000000000$ (first-order: 11 contexts; second-order: 15 contexts); all totals close to 1: True (full list in the appendix output). \textbf{Q8.} A total of $0.87$ would mean probability mass is missing: the implementation has a bug, e.g.\ dividing by the wrong total (such as a total that includes another context or is computed before all counts are added), dropping the \texttt{<END>} transition, or filtering out some tokens.

\section*{Part VIII: Next-word prediction (Question 9)}
$\arg\max$ for the first-order model: after \emph{the}: cat (tie with dog at 0.25; tie broken by first insertion); after \emph{cat}: sat (0.667); after \emph{dog}: sat (0.667); after \emph{sat}: on; after \emph{ran}: to; after \emph{on}: the; after \emph{to}: the.
\textbf{Q9.} Often yes but not always: after \emph{the}, a human would not prefer \emph{cat} over \emph{dog}; the model's choice is an artefact of a tie and of only six sentences. The model reflects the frequencies in its data, not meaning or world knowledge; human expectations come from broad experience, so the two coincide only where the data agree with them.

\section*{Part IX--X: Generation (Question 10)}
20 first-order sampled sentences (seed 0), e.g.\ ``the dog sat on the mat'', ``the rug'', ``the park'', ``the cat sat on the park'', and also ``the cat ran to the dog sat on the cat sat on the rug'' (full list in the appendix). Greedy vs sampling (5 each): greedy produced the same sentence five times, a loop ``the cat sat on the cat sat on the \dots'' that never reaches \texttt{<END>} and is cut off by \texttt{max\_len}=30; sampling gave varied sentences (``the cat ran to the rug'', ``the cat sat on the rug'', ``the park'', \dots).
\textbf{Q10.} Sampling produces more variation, because it follows the whole distribution, while greedy always chooses the same next word for a given context and so is deterministic (here it also gets trapped in a cycle in the first-order model).

\section*{Part XI--XII: Second-order model (Question 11)}
Prompt 2 in the appendix. Before accepting the code, what should change was specified: the context becomes a pair $(X_{t-2},X_{t-1})$, the CPT is indexed by pairs, and sentences need \emph{two} \texttt{<START>} tokens so that the first word has a full context. The generalised class \texttt{NGramLM(order=2)} does this.
Selected second-order CPTs: $P(\cdot\mid\texttt{<START>},\text{the})$: cat 0.5, dog 0.5; $P(\cdot\mid\text{the},\text{cat})$: sat 0.667, ran 0.333; $P(\cdot\mid\text{on},\text{the})$: mat 0.5, rug 0.5; $P(\cdot\mid\text{to},\text{the})$: park 1.0.
\textbf{Q11.} (1) Graph: $X_t$ has two parents, $X_{t-2}\to X_t\leftarrow X_{t-1}$, instead of one. (2) CPT: rows indexed by pairs, up to $|V|^2$ rows of size $|V|$ instead of $|V|$ rows. (3) Context: two previous words, so it can separate ``on the $\to$ mat/rug'' from ``to the $\to$ park'' (first-order ``the'' mixes all five continuations). (4) Data: many more contexts so many more observations are needed; most pair contexts are never seen.

\section*{Part XIII: Comparison (Question 12)}
\begin{center}\begin{tabular}{lcc}\toprule Measure & First-order & Second-order\\\midrule
Non-zero CPT entries & 17 & 19\\
Observed contexts & 11 & 15\\
Possible contexts (incl.\ \texttt{<START>}) & 11 & 121\\
Unobserved (zero-probability) contexts & 0 & 106 (88\%)\\
Full CPT size (contexts $\times|V|$) & 121 & 1331\\
Distinct sentences in 5000 samples & 482 & 6\\
Novel sentences (not in training) & 476 & 0\\\bottomrule\end{tabular}\end{center}
Examples. First-order (sampled): ``the park'', ``the cat sat on the park'', ``the dog sat on the dog ran to the park'': varied but often ungrammatical or incoherent. Second-order (sampled): all 20 samples are grammatical (``the cat ran to the park''), but they are only the six training sentences: it reproduces the data with no novelty. Sentence probabilities: ``the dog ran to the mat'' has $P=0.0139$ under the first-order model but $0$ under the second-order model (context ``to the'' was only followed by park); ``the cat sat on the mat'' gets $0.0278$ vs $0.1667$.
\textbf{Q12.} More context separates situations that a short context merges (e.g.\ ``on the'' vs ``to the''), so predictions become sharper and the generated text more coherent. But the CPT grows as $|V|^{n}\cdot|V|$ (121 to 1331 entries here) while the data stay fixed, so most contexts are never observed (88\%) or seen once, the estimates are unreliable, unseen contexts get no prediction at all, and the model memorises the data. Context helps bias but hurts variance.

\section*{Part XIV--XV: Questions 13--14}
\textbf{Q13.} Approach B states the intended model, so (i) behaviour is specified in advance and can be checked against it, (ii) I must understand the representation (CPT of counts) to state it, (iii) the output can be validated against the specification, e.g.\ by examining where counts and probabilities are stored, (iv) probabilistic invariants such as $\sum_vP(v\mid w)=1$ and the no-unseen-context behaviour can be tested, and (v) I can separate the model (the conditional distribution $P(X_t\mid X_{t-1})$) from its implementation (dictionaries, sampling code). With Approach A the LLM could pick a neural network, smoothing, or something else, and I would have nothing to test against.

\textbf{Q14.} (a) \emph{Factorisation of the joint:} $P(x_1..x_T)=\prod_tP(x_t\mid\text{parents})$, which gave the sentence probabilities computed above. (b) \emph{Independence assumptions:} the graph makes the Markov assumption explicit ($X_t\perp X_{1:t-2}\mid X_{t-1}$) and shows exactly what the second-order model relaxes. (c) \emph{Effect of context:} adding a parent adds a dimension to the CPT, so the table size grows exponentially in the number of parents (121 to 1331). (d) \emph{Principled generation:} ancestral sampling along the graph, sample $X_1$, then each child given its sampled parents. (e) \emph{Testing:} each CPT row must be a distribution, so the normalisation check tests that the implementation matches the probabilistic specification. The same view carries over to modern LMs: only the representation of $P(X_t\mid X_{<t})$ changes (a neural network instead of a table).

\section*{Reflection on the use of the LLM}
The LLM wrote the class, the CPT construction, sampling, generation and the tests, which would have taken much longer by hand; The model was specified first (Approach B) and tested against the specification. \textbf{Code inspection findings (no corrections were needed):} (1) in \texttt{generate()}, a greedy loop with no length limit never terminates in the first-order model, because ``the cat sat on the cat sat\dots'' is a cycle that never reaches \texttt{<END>} (observed output above), so the \texttt{max\_len} cap is essential; (2) handling of unseen contexts returns \texttt{None} rather than raising \texttt{KeyError}; (3) the second-order extension needs \emph{two} \texttt{<START>} tokens, otherwise the first word has no context and the $\sum=1$ test would not cover the first word correctly. It was also checked that sampling uses \texttt{random.choices} with probabilities (not argmax), and the normalisation test was run on both models.

\section*{Appendix A: Prompts}\VerbatimInput{prompts.txt}
\section*{Appendix B: Program output}\VerbatimInput{results.txt}
\section*{Appendix C: Code (\texttt{ngram\_lm.py})}\VerbatimInput{ngram_lm.py}
\end{document}
